\section{Proposed Joint Quantization and Detection Framework}

\subsection{AP-Side Triple-Head Policy Network}
To enable task-oriented dynamic quantization, each AP $l$ is equipped with a lightweight triple-head Policy Network, denoted as $\pi_\phi$. Its role is to independently infer the optimal bit-widths $b_{l,i}$ for the three components $i \in \{y, H, s\}$. 

The policy network is modeled as a 3-layer Multi-Layer Perceptron (MLP) with a hidden dimension of 64 and ReLU activation functions. The input is a locally extracted statistical feature vector $\mathbf{a}_l \in \mathbb{R}^{3}$, which comprises the local received signal power $\|\mathbf{y}_l\|^2/N$, the Frobenius norm of the estimated channel $\|\hat{\mathbf{H}}_l\|_F^2$, and the local empirical Signal-to-Interference-plus-Noise Ratio (SINR).

The output of each head is a probability logits vector for the categorical distribution over the available bit-widths $\mathcal{B} = \{0, 4, 8, 12, 16\}$. To facilitate end-to-end differentiable backpropagation through discrete categorical sampling, we utilize the Gumbel-Softmax relaxation. The continuous approximation of the categorical selection weight for the $j$-th bit-width option is given by:
\begin{equation}
    w_{l,i}^{(j)} = \frac{\exp((\log \pi_\phi(\mathbf{a}_l)_{i,j} + g_{l,i,j})/\tau)}{\sum_{m=1}^{|\mathcal{B}|} \exp((\log \pi_\phi(\mathbf{a}_l)_{i,m} + g_{l,i,m})/\tau)},
\end{equation}
where $g_{l,i,j}$ are i.i.d. noise samples drawn from the standard Gumbel distribution $\text{Gumbel}(0,1)$, and $\tau > 0$ is the annealing temperature parameter. The effective differentiable bit-width utilized in the continuous LSQ penalty during training is computed as $b_{l,i} = \sum_{j=1}^{|\mathcal{B}|} w_{l,i}^{(j)} \mathcal{B}_j$.

\subsection{CPU-Side Detection with Bit-Aware Attention}
Upon receiving the quantized features, the CPU employs a hybrid architecture comprising a Mean-Field Graph Neural Network (GNN) for initial spatial interference cancellation, followed by a Transformer encoder for refined multi-user detection.

\subsubsection{Mean-Field GNN Layer}
The CF-MIMO network is modeled as a bipartite graph connecting APs and UEs. For a user node $k$ at iteration $t$, the GNN aggregates spatial messages from all connected APs. Crucially, the bit-width $b_{l,y}$ is concatenated with the quantized physical features to inform the network of the varying fidelity of the messages:
\begin{align}
    \mathbf{m}_{lk}^{(t)} &= \text{MLP}_{msg}\left( \left[ \mathbf{h}_{lk}^q \parallel \mathbf{y}_l^q \parallel b_{l,y} \right] \right), \\
    \mathbf{z}_{k}^{(t+1)} &= \text{MLP}_{update}\left( \mathbf{z}_k^{(t)}, \frac{1}{L} \sum_{l=1}^L \mathbf{m}_{lk}^{(t)} \right),
\end{align}
where $\parallel$ denotes concatenation, $\mathbf{z}_k^{(t)}$ is the hidden state of user $k$, and the aggregation function is implemented as a spatial MEAN pooling over $L$ APs. Explicitly passing the bit-width enables the GNN to automatically down-weight noisy messages originating from heavily quantized or dropped ($b=0$) AP links.

\subsubsection{Bit-Aware Transformer}
To capture complex inter-user interference patterns over the aggregated features $\mathbf{Z} \in \mathbb{R}^{K \times d_{model}}$, we deploy a Transformer encoder equipped with a novel \textit{Bit-Aware Multi-Head Attention (BA-MHA)}. Standard attention relies solely on the feature content via query $\mathbf{Q} = \mathbf{Z}\mathbf{W}_Q$ and key $\mathbf{K} = \mathbf{Z}\mathbf{W}_K$, where $\mathbf{W}_Q, \mathbf{W}_K \in \mathbb{R}^{d_{model} \times d_k}$ are learnable linear projection matrices.

We introduce an explicit topological bias based on the overall quantization resolution of the contributing APs. Let $\mathbf{e}_{l}$ be the one-hot encoding of the bit-widths assigned to AP $l$. The modified attention coefficient from user $k$ to an AP representation $l$ is expressed as:
\begin{equation}
    \alpha_{kl} = \frac{\exp\left( \frac{\mathbf{q}_k^T \mathbf{k}_l}{\sqrt{d_k}} + \text{MLP}_{bit}(\mathbf{e}_l) \right)}{\sum_{l'} \exp\left( \frac{\mathbf{q}_k^T \mathbf{k}_{l'}}{\sqrt{d_k}} + \text{MLP}_{bit}(\mathbf{e}_{l'}) \right)}.
\end{equation}
This architecture enforces a form of task-oriented spatial filtering, dynamically focusing the CPU's processing capacity on high-precision spatial streams.

\subsection{Training and Implementation Details}
To guarantee full reproducibility, we detail the end-to-end training strategy. The dataset is dynamically generated during training based on the 3GPP UMi channel model, simulating random large-scale fading topologies in a $1\times 1$ km$^2$ area. We utilize the Adam optimizer across 200 total epochs, grouped into batches of size 256. 

The training is strictly divided into two phases. In Phase 1 (50 epochs), we train the hybrid GNN-Transformer detector under fixed high-resolution uniform quantization (learning rate $10^{-3}$) to initialize a stable detection backbone. In Phase 2 (150 epochs), we activate the Policy Network and LSQ modules to jointly learn the dynamic bit allocation and detection (learning rate $5\times 10^{-4}$). The Gumbel-Softmax temperature $\tau$ is linearly annealed from $1.0$ down to $0.1$ to encourage hard, discrete bit decisions. The Lagrangian penalty factor $\lambda$ is tuned precisely between $[0.0, 0.2]$ to establish the Pareto frontier of the detection BER versus average fronthaul bit consumption.